\documentclass[10pt,a4paper]{amsart}
\usepackage{amsmath,amssymb,amsthm,graphicx}
% Retained preamble for the Stacks Project chapter {CH} (shared project preamble).
% The chapter file pinned-original-source.tex only carries \input{{preamble}} on its first line,
% so no preamble text lives inside the chapter itself.  The Stacks Project uses one single shared
% preamble for every chapter: the snapshot stages it three times (sources/stacks-chow/preamble.tex,
% sources/stacks-etale-cohomology/preamble.tex and focused-seeds/AGP/sources/stacks-snapshot-20261011/
% preamble.tex) and those copies are identical in size.  The block below is copied mechanically,
% line for line, from sources/stacks-chow/preamble.tex lines 14-18 and 154-217: the Xy-pic setup and
% the native theorem declarations together with the mathematical macros.  Nothing was added, removed
% or redefined; the wrapper class and the amsmath/amssymb/amsthm/graphicx base packages are supplied
% by the assembler.
% For commutative diagrams we use Xy-pic
% Chapter: stacks-cohomology.tex -- Cohomology of Sheaves.
\usepackage[all]{xy}

% We use 2cell for 2-commutative diagrams.
\xyoption{2cell}

% Theorem environments.
%
\theoremstyle{plain}
\newtheorem{theorem}[subsection]{Theorem}
\newtheorem{proposition}[subsection]{Proposition}
\newtheorem{lemma}[subsection]{Lemma}

\theoremstyle{definition}
\newtheorem{definition}[subsection]{Definition}
\newtheorem{example}[subsection]{Example}
\newtheorem{exercise}[subsection]{Exercise}
\newtheorem{situation}[subsection]{Situation}

\theoremstyle{remark}
\newtheorem{remark}[subsection]{Remark}
\newtheorem{remarks}[subsection]{Remarks}

\numberwithin{equation}{subsection}

% Macros
%
\def\lim{\mathop{\mathrm{lim}}\nolimits}
\def\colim{\mathop{\mathrm{colim}}\nolimits}
\def\Spec{\mathop{\mathrm{Spec}}}
\def\Hom{\mathop{\mathrm{Hom}}\nolimits}
\def\Ext{\mathop{\mathrm{Ext}}\nolimits}
\def\SheafHom{\mathop{\mathcal{H}\!\mathit{om}}\nolimits}
\def\SheafExt{\mathop{\mathcal{E}\!\mathit{xt}}\nolimits}
\def\Sch{\mathit{Sch}}
\def\Mor{\mathop{\mathrm{Mor}}\nolimits}
\def\Ob{\mathop{\mathrm{Ob}}\nolimits}
\def\Sh{\mathop{\mathit{Sh}}\nolimits}
\def\NL{\mathop{N\!L}\nolimits}
\def\CH{\mathop{\mathrm{CH}}\nolimits}
\def\proetale{{pro\text{-}\acute{e}tale}}
\def\etale{{\acute{e}tale}}
\def\QCoh{\mathit{QCoh}}
\def\Ker{\mathop{\mathrm{Ker}}}
\def\Im{\mathop{\mathrm{Im}}}
\def\Coker{\mathop{\mathrm{Coker}}}
\def\Coim{\mathop{\mathrm{Coim}}}

% Boxtimes
%
\DeclareMathSymbol{\boxtimes}{\mathbin}{AMSa}{"02}

%
% Macros for moduli stacks/spaces
%
\def\QCohstack{\mathcal{QC}\!\mathit{oh}}
\def\Cohstack{\mathcal{C}\!\mathit{oh}}
\def\Spacesstack{\mathcal{S}\!\mathit{paces}}
\def\Quotfunctor{\mathrm{Quot}}
\def\Hilbfunctor{\mathrm{Hilb}}
\def\Curvesstack{\mathcal{C}\!\mathit{urves}}
\def\Polarizedstack{\mathcal{P}\!\mathit{olarized}}
\def\Complexesstack{\mathcal{C}\!\mathit{omplexes}}
% \Pic is the operator that assigns to X its picard group, usage \Pic(X)
% \Picardstack_{X/B} denotes the Picard stack of X over B
% \Picardfunctor_{X/B} denotes the Picard functor of X over B
\def\Pic{\mathop{\mathrm{Pic}}\nolimits}
\def\Picardstack{\mathcal{P}\!\mathit{ic}}
\def\Picardfunctor{\mathrm{Pic}}
\def\Deformationcategory{\mathcal{D}\!\mathit{ef}}
\usepackage{cleveref}
\crefname{originaldefinition}{original definition}{original definitions}
\Crefname{originaldefinition}{Original definition}{Original definitions}
\makeatletter
\newlabel{algebra-definition-derivation}{{00RN}{source}{}{sourcepointer-1}{}}
\newlabel{algebra-definition-derivation@cref}{{[originaldefinition][0][]00RN}{[1][0][]source}{}{}{}}
\makeatother
\crefname{originallemma}{original lemma}{original lemmas}
\Crefname{originallemma}{Original lemma}{Original lemmas}
\makeatletter
\newlabel{lemma-cech-h1}{{0B8R}{source}{}{sourcepointer-2}{}}
\newlabel{lemma-cech-h1@cref}{{[originallemma][0][]0B8R}{[1][0][]source}{}{}{}}
\makeatother
\title{The boundary map of a square zero ideal extension of sheaves of rings is a derivation}
\author{The Stacks Project Authors}
\date{stacks-project:cohomology.tex:local-snapshot-20261011; GNU FDL 1.2 or later}
\begin{document}
\maketitle
\noindent\textit{Editorial note.} All delivered mathematical text is the author's own, copied from the pinned Stacks Project chapter file cohomology.tex; no statement or proof step was rewritten, reordered or removed.\par\medskip
\noindent\textit{Editorial note.} The proof is complete as written and ends with the displayed computation of the cocycle of the product; no step is referred to the reader.\par\medskip
\noindent\textit{Editorial note.} The counted spans carry 2 navigation stubs and no bibliography citation.\par\medskip
\noindent\textit{Editorial note.} The preamble is the chapter's retained minimal preamble (sources/stacks-cohomology/retained-preamble.tex lines 1-81), a verbatim excerpt of the native preamble.\par\medskip
\noindent\textit{Editorial note.} The label lemma-boundary-derivation-over-cup-product is defined only in cohomology.tex, so no cross-chapter restatement had to be weighed for this card.\par\medskip
\noindent\textit{Allowed standard prerequisites (editorial).}\begin{itemize}
\item Boundary map of a short exact sequence of sheaves in cohomology. Same chapter, the long exact cohomology sequence of a short exact sequence of sheaves, used at proof lines 5746-5747. Supplies the map whose derivation property is being proved.
\item Description of H\textasciicircum{}1 by Cech cocycles for a covering. Same chapter, Lemma lemma-cech-h1, invoked at proof lines 5766-5767. Supplies that the first Cech cohomology group of the covering embeds in M, so that the boundary classes can be represented by cocycles.
\item Notion of a derivation of a ring into a module. Algebra, Definition algebra-definition-derivation, named in the statement at line 5751. Supplies the two axioms (additivity and the Leibniz rule) that the boundary map has to satisfy.
\end{itemize}\medskip
\ifdefined\hypertarget\hypertarget{sourcepointer-1}{}\fi
\noindent\textit{Source reference.} Original definition 00RN: sources/stacks-chow/tags.txt, source lines 1009--1009. Pointer outside the counted span, invoked once for context; the load-bearing argument of the counted statement is the author's proof reproduced in full in the proof span.\par\medskip
\ifdefined\hypertarget\hypertarget{sourcepointer-2}{}\fi
\noindent\textit{Source reference.} Original lemma 0B8R: sources/stacks-cohomology/pinned-original-source.tex, source lines 1486--1486. Pointer outside the counted span, invoked once for context; the load-bearing argument of the counted statement is the author's proof reproduced in full in the proof span.\par\medskip
\begin{lemma}
\label{lemma-boundary-derivation-over-cup-product}
Let $X$ be a topological space. Let $\mathcal{O}' \to \mathcal{O}$ be a
surjection of sheaves of rings whose kernel $\mathcal{I} \subset \mathcal{O}'$
has square zero. Then $M = H^1(X, \mathcal{I})$ is a
$R = H^0(X, \mathcal{O})$-module and the boundary map
$\partial : R \to M$ associated to the short exact sequence
$$
0 \to \mathcal{I} \to \mathcal{O}' \to \mathcal{O} \to 0
$$
is a derivation (Algebra, Definition \ref{algebra-definition-derivation}).
\end{lemma}

\begin{proof}
The map $\mathcal{O}' \to \SheafHom(\mathcal{I}, \mathcal{I})$
factors through $\mathcal{O}$ as $\mathcal{I} \cdot \mathcal{I} = 0$
by assumption. Hence $\mathcal{I}$ is a sheaf of $\mathcal{O}$-modules
and this defines the $R$-module structure on $M$.
The boundary map is additive hence it suffices to prove
the Leibniz rule. Let $f \in R$. Choose an open covering
$\mathcal{U} : X = \bigcup U_i$ such that there exist
$f_i \in \mathcal{O}'(U_i)$ lifting $f|_{U_i} \in \mathcal{O}(U_i)$.
Observe that $f_i - f_j$ is an element of $\mathcal{I}(U_i \cap U_j)$.
Then $\partial(f)$ corresponds to the {\v C}ech cohomology class of
the $1$-cocycle $\alpha$ with $\alpha_{i_0i_1} = f_{i_0} - f_{i_1}$.
(Observe that by Lemma \ref{lemma-cech-h1} the first {\v C}ech cohomology
group with respect to $\mathcal{U}$ is a submodule of $M$.)
Next, let $g \in R$ be a second element and assume (after possibly
refining the open covering) that $g_i \in \mathcal{O}'(U_i)$ lifts
$g|_{U_i} \in \mathcal{O}(U_i)$. Then we see that
$\partial(g)$ is given by the cocycle $\beta$ with
$\beta_{i_0i_1} = g_{i_0} - g_{i_1}$. Since $f_ig_i \in \mathcal{O}'(U_i)$
lifts $fg|_{U_i}$ we see that
$\partial(fg)$ is given by the cocycle $\gamma$ with
$$
\gamma_{i_0i_1} = f_{i_0}g_{i_0} - f_{i_1}g_{i_1} =
(f_{i_0} - f_{i_1})g_{i_0} + f_{i_1}(g_{i_0} - g_{i_1}) =
\alpha_{i_0i_1}g + f\beta_{i_0i_1}
$$
by our definition of the $\mathcal{O}$-module structure on $\mathcal{I}$.
This proves the Leibniz rule and the proof is complete.
\end{proof}

\end{document}
